\section{Bounds and Limitations}\label{sec:bounds}

Describing the Kasai construction inside a product group buys more than notational convenience: because the lift group factors, so do the objects the code parameters depend on, and the rate, distance and girth of a GALA code can be bounded directly in terms of the four design parameters $L$, $J$, $k$ and $m$. This section collects those bounds. They are limitations as much as guarantees --- read in the contrapositive, they say which regions of the parameter space cannot contain a good code, and it is precisely the boundaries they draw that motivate the polynomial and loose-active generalizations of Definitions~\ref{def:pg-code} and~\ref{def:loose-orthogonality}.

Two scoping conventions are used throughout. First, unless stated otherwise the results of this section assume a \emph{monomial} lift, in which every entry of $\mathcal F$ and $\mathcal G$ is a single group element rather than a sum; this is exactly the regime in which the stabilizer weight is $w = L$ and every data qubit meets exactly $J$ checks of each type. Both facts are used constantly below, and both fail for polynomial lifts. Second, the letter $k$ does double duty in this paper, as the number of logical qubits of an $[[n,k,d]]$ code and as the degree of the top permutation group $H_k$ acting on $[k]$, so that $n = Lkm$. We keep the established notation and let context disambiguate, noting only that the two never meet inside a displayed formula below: the rate bound~\eqref{eq:mono-rate} is written in terms of $n$ precisely so that it does not. Note also that $k$ is the \emph{degree} of $H_k$ and not its order, which for the tops used here is larger --- $|S_3| = 6$ at $k=3$, $|S_4| = 24$ at $k=4$.

\subsection{Bounds on Monomial Direct Product Constructions}\label{ssec:mono-bounds}

We begin with the bounds that follow from each generator being a permutation matrix.

\begin{proposition}[Rate Bound]\label{prop:rate-bound}
Let $\mathcal Q\in \mathrm{GALA}_{L,J}(H_k\star C_m)$ be an $[[n,k,d]]$ monomial GALA code. Then
\begin{equation}\label{eq:mono-rate}
    \frac{k}{n} \geq 1-\frac{2J}{L} + \frac{2(J-1)}{n}.
\end{equation}
In particular the rate exceeds $1-2J/L$ strictly whenever $J\geq2$, and is at least $\tfrac12$ whenever $J\leq \tfrac L4$.
\end{proposition}
\begin{proof}
    All rows sum up to $\textbf{1}^{1\times n}$, contributing to the $J-1$ rank deficiencies.     
\end{proof}


\begin{proposition}(Distance Bounds from J,L)\label{prop:distance-bound}
Let $\mathcal{Q}\in \mathrm{GALA}_{L,J}(H_k\star C_m)$ be an $[[n,k,d]]$ GALA code. Then,

\begin{align}
     J\leq 2        &\implies d\leq \frac12\gth{\mathcal{Q}}\label{eq:j<=2}, \\
     J>\frac{L}{4}  &\implies d\leq L,\label{eq:j>=L/4}\\
     L<12&\implies d\leq L.\label{eq:j<12}
\end{align}

\end{proposition}

\begin{proof}
    When $J\leq 2$, each data qubit $q$ belongs to exactly two stabilizers $c, c'$, and any Tanner graph cycles containing $q$ must contain both $c,c'$. It follows that every $t$-cycle is a wt-$\frac t2$ logical. When $J> \frac L4$, the parent matrix $\hat{H}_X, \hat{H}_Z$ satisfy the full CSS orthogonality $\implies d\leq \wt{\hat H} = L$. Finally, the first two conditions implies the third. 
\end{proof}


The product structure naturally preserves many properties of its factors, giving us another class of bounds. We can analyze these bounds by first defining the quotient codes:

\begin{definition}[Quotient Codes]\label{def:quotient}
Suppose $\mathcal{Q}\in \mathrm{GALA}_{L,J}(H_k\times C_m)$ is a monomial DPG code with generators $\mathcal{F} = \{(F_i, f_i)\}_{i\in[\frac L2]}$ and $\mathcal{G} = \{(G_i,g_i)\}_{i\in [\frac L2]}$, so that the lifted blocks are the Kronecker products $F_i\otimes f_i$ and $G_i\otimes g_i$ of Definition~\ref{def:direct-product}. Discarding one tensor factor of every generator leaves a smaller GALA code over the surviving factor: the \emph{top quotient code} $\mathcal{Q}^H = \mathrm{GALA}_{L,J}(\{F_i\}, \{G_i\})$ on $Lk$ qubits and the \emph{bottom quotient code} $\mathcal{Q}_C = \mathrm{GALA}_{L,J}(\{f_i\}, \{g_i\})$ on $Lm$ qubits, both with the same $L$, $J$ and the same block-circulant row order as $\mathcal Q$. If moreover $C_m =\mathbb{Z}_{p_1}\times \cdots \times \mathbb{Z}_{p_q}$ with $\gcd(p_a, p_b) = 1$ for $a\neq b$, then each bottom generator factors as $f_i = f^{(1)}_i\otimes\cdots\otimes f^{(q)}_i$ and $g_i = g^{(1)}_i\otimes\cdots\otimes g^{(q)}_i$ with $f^{(a)}_i,g^{(a)}_i\in\mathbb Z_{p_a}$, and the bottom quotient factors further into the $q$ codes $\mathcal{Q}_{C_a} = \mathrm{GALA}_{L,J}(\{f^{(a)}_i\}, \{g^{(a)}_i\})$ on $Lp_a$ qubits each.
\end{definition}

Each of these quotients is obtained by summing over the orbits of a normal subgroup of the lift group. Concretely, let $\pr_{C} = I_{k}\otimes \mathbf{1}^{1\times m}$ and $\pr^{H} = \mathbf{1}^{1\times k}\otimes I_m$ be the two orbit-sum matrices of Section~\ref{ssec:notations}, so that $\pr_C$ collapses the $C_m$ coordinate and $\pr^H$ collapses the $H_k$ coordinate: $\pr_C(h,c) = h$ and $\pr^H(h,c) = c$ for all $(h,c)\in H_k\times C_m$, the latter because in a direct product the $C_m$ component is the same in every block. Applying $\pr_C$ blockwise to $\hat H_X$ therefore returns the parent matrix of $\mathcal Q^H$, and $\pr^H$ that of $\mathcal Q_C$.

Transporting a logical operator the other way, from a quotient back up to $\mathcal Q$, is the operation that converts these factorizations into distance bounds.

\begin{lemma}[Logical Inflation]\label{lem:inflation}
Let $\mathcal{Q}\in \mathrm{GALA}_{L,J}(H_k\times C_m)$ be a monomial DPG code and let $\mathcal Q_N$ be the quotient obtained by collapsing a tensor factor of size $\mu$, so that the qubits of $\mathcal Q$ are indexed by $[n_N]\times[\mu]$ with $[n_N]$ indexing the qubits of $\mathcal Q_N$. Write $\Pi\in\mathbb F_2^{n\times n_N}$ for the corresponding orbit-sum matrix, $\Pi_{(a,x),b} = \delta_{ab}$, and define the \emph{inflation} of a row vector $\ell\in\mathbb F_2^{1\times n_N}$ to be $\ell' := \ell\,\Pi^T$, i.e.\ $\ell'_{(a,x)} = \ell_a$. Then:
\begin{enumerate}
    \item $\wt{\ell'} = \mu\cdot\wt{\ell}$;
    \item $\ell\in\ns{H_Z(\mathcal Q_N)} \iff \ell'\in\ns{H_Z(\mathcal Q)}$,
    \item $\ell\in\rs{H_X(\mathcal Q_N)} \implies \ell'\in\rs{H_X(\mathcal Q)}$, 
\end{enumerate}
and likewise with $X$ and $Z$ exchanged. 
\end{lemma}

\begin{proof}
Index the qubits of $\mathcal Q$ by triples $(a,i,x)$ with $a\in[L]$ the block column, $i$ the surviving coordinate and $x\in[\mu]$ the collapsed one; by Definition~\ref{def:direct-product} the lifted block in block position $(r,a)$ of $\hat H_X$ is a Kronecker product $A_{a-r}\otimes B_{a-r}$ with $A$ the surviving and $B$ the collapsed factor, both permutation matrices because the lift is monomial. Claim (1) is immediate, every qubit of $\mathcal Q_N$ being replaced by exactly $\mu$ qubits of $\mathcal Q$.

For (2), compute the syndrome of $\ell'$ at the check indexed by $(r,i',x')$:
\[
    \big(H_X\ell'^T\big)_{(r,i',x')}
    = \sum_{a,i,x}[A_{a-r}]_{i',i}\,[B_{a-r}]_{x',x}\,\ell_{(a,i)}
    = \sum_{a,i}[A_{a-r}]_{i',i}\,\ell_{(a,i)}
    = \big(H_X(\mathcal Q_N)\,\ell^T\big)_{(r,i')} ,
\]
where the middle step used $\sum_x [B_{a-r}]_{x',x} = 1$, valid because $B_{a-r}$ is a permutation matrix and hence has exactly one $1$ in each row. So the syndrome of $\ell'$ is the syndrome of $\ell$, replicated $\mu$ times, and one vanishes iff the other does. The same computation applies verbatim to $H_Z$, whose blocks $G^T_{r-a} = A^T_{r-a}\otimes B^T_{r-a}$ have the same Kronecker form.

For (3), suppose $\ell = uH_X(\mathcal Q_N)$ and set $u'_{(r,i',x')} := u_{(r,i')}$. Then
\[
    \big(u'H_X\big)_{(a,i,x)} = \sum_{r,i',x'}u_{(r,i')}[A_{a-r}]_{i',i}[B_{a-r}]_{x',x}
    = \sum_{r,i'}u_{(r,i')}[A_{a-r}]_{i',i} = \ell_{(a,i)} = \ell'_{(a,i,x)},
\]
using $\sum_{x'}[B_{a-r}]_{x',x} = 1$, this time down a column. Hence $\ell' = u'H_X\in\rs{H_X}$. Together (2) and (3) say that inflation carries $\ns{H_Z(\mathcal Q_N)}$ into $\ns{H_Z(\mathcal Q)}$ and $\rs{H_X(\mathcal Q_N)}$ into $\rs{H_X(\mathcal Q)}$, so it descends to the quotients and defines the asserted map on $\mathcal L_X = \ns{H_Z}/\rs{H_X}$.
\end{proof}


\begin{corollary}(Distance Inheritance from Quotients)\label{cor:quotient-distance}
Let $\mathcal{Q}\in \mathrm{GALA}_{L,J}(H_k\times C_m)$ be an $[[n,k,d]]$ direct-product GALA code. Let $\mathcal{Q}^H$, $\mathcal{Q}_C$ be the top and bottom quotient codes. Then,

\begin{equation}
    d_\mathcal{Q}\leq \min[m\cdot d_{\mathcal{Q}^H},\, k\cdot d_{\mathcal{Q}_C}].
\end{equation}

In general, whenever $C_m =  \mathbb  Z_p\times \mathbb Z_q$, we also have that
\begin{equation}
    d_\mathcal{Q}\leq k\cdot\min[ q\cdot  d_{\mathcal{Q}_{C_p}}, \, p\cdot d_{\mathcal{Q}_{C_q}}].
\end{equation}
\end{corollary}
\begin{proof}
    Follows from  Lemma~\ref{lem:inflation} by setting $N = H_k, C_m, C_p, C_q$.
\end{proof}

\begin{corollary}(Distance Bounds for QuEra Kasai)
    Let $\mathcal{Q}\in \mathrm{GALA}_{L,J}(H_k\times C_m)$ be an $[[n,k,d]]$ direct-product GALA code. We have 

    \begin{equation}
        d\leq kL.
    \end{equation}
    Furthermore, if $H_k = S_3$, we have that 
    \begin{equation}
        d\leq \min [3L,\, 2m ].
    \end{equation}
\end{corollary}

\begin{proof}
    Bottom code has distance $\leq L$ since it's fully orthogonal (non-abelian factors are projected away). Top code has distance $2$ for $S_3$ since there can be at most $4$ non-identity elements. We either have full orthgonality, which forces $d\leq L$, or we have two identical columns of identities in $\mathcal{Q}^H$, which implies $d_{\mathcal{Q}^H}\leq 2$.
\end{proof}


\begin{lemma}[Girth Inheritance from Quotients]\label{lem:girth-inheritance}
Let $\mathcal{Q}\in \mathrm{GALA}_{L,J}(H_k\times C_m)$ be a monomial DPG code, and let $\mathcal Q_N$ be any quotient code of Definition~\ref{def:quotient}. Then
\begin{equation}\label{eq:girth-inheritance}
    \gth{\mathcal{Q}}\geq \gth{\mathcal{Q}_{N}}.
\end{equation}
\end{lemma}

\begin{proof}
Write the qubits of $\mathcal Q$ as pairs $(a,x)$ with $a$ a qubit of $\mathcal Q_N$ and $x\in[\mu]$ the collapsed coordinate, and likewise its checks as $(c,x')$, and let $\pi$ be the map forgetting the second coordinate. As in the proof of Lemma~\ref{lem:inflation}, the block of $\hat H_X$ in position $(r,a)$ is a Kronecker product $A_{a-r}\otimes B_{a-r}$ of permutation matrices, so the qubit $(a,x)$ is incident to the check $(c,x')$ exactly when $a$ is incident to $c$ in $\mathcal T(\mathcal Q_N)$ \emph{and} $x = B(x')$ for the single permutation $B$ attached to that block. Hence for each fixed check $(c,x')$ and each neighbor $a$ of $c$ downstairs there is exactly one $x$ with $(a,x)$ a neighbor of $(c,x')$, and symmetrically for a fixed qubit: $\pi$ restricts to a bijection between the neighborhood of any vertex of $\mathcal T(\mathcal Q)$ and the neighborhood of its image. That is, $\pi$ is a graph covering.

A covering map sends a cycle of length $t$ in $\mathcal T(\mathcal Q)$ to a closed walk of length $t$ in $\mathcal T(\mathcal Q_N)$ which is non-backtracking, because $\pi$ is injective on neighborhoods; a non-backtracking closed walk of length $t$ contains a cycle of length at most $t$. Therefore every cycle upstairs has length at least $\gth{\mathcal Q_N}$, which is~\eqref{eq:girth-inheritance}. The same argument applies in the $Z$ sector, and $\gth{\mathcal Q}$ is the minimum of the two.
\end{proof}

\begin{corollary}[Girth Bound on $m$]\label{cor:girth-m}
Fix $L$, $J$ and a top $H_k$ admitting an active set $\Gamma$, and let $C_m = \mathbb Z_m$ be cyclic. If $m> 2(J-1)L$ then there is a choice of bottom generators for which the resulting $\mathcal{Q}^*\in \mathrm{GALA}_{L,J}(H_k\times C_m)$ has $\gth{\mathcal{Q}^*}\geq 6$.
\end{corollary}

\begin{proof}
By Lemma~\ref{lemma:girth} it suffices to make every off-diagonal entry of $H_XH_X^T$ and of $H_ZH_Z^T$ at most $1$ over $\mathbb Z$. Write $\ell = \tfrac L2$, let $x$ denote the generator of $\mathbb Z_m$, and let the lifts be $F_u = F^H_u\otimes x^{f_u}$ and $G_u = G^H_u\otimes x^{g_u}$. The $(i,i')$ block of $H_XH_X^T$ depends only on $\delta = i'-i$ and equals
\[
    B_\delta = \sum_{u\in[\ell]}\Big(F^H_u F^{H,T}_{u-\delta}\otimes x^{\,f_u-f_{u-\delta}}\Big)
             + \sum_{u\in[\ell]}\Big(G^H_u G^{H,T}_{u-\delta}\otimes x^{\,g_u-g_{u-\delta}}\Big),
\]
a sum of $L$ permutation matrices. At $\delta = 0$ every summand is the identity, so $B_0 = L\cdot I$ is diagonal and contributes no off-diagonal entry. For $\delta\neq0$, two summands can share a non-zero position only if they carry the same power of $S$; so it suffices that, for each of the $J-1$ offsets $\delta = 1,\dots,J-1$ realizable as $i'-i$ with $i,i'\in[J]$, the $L$ differences
\[
    D_\delta \;=\; \big(f_u-f_{u-\delta}\big)_{u\in[\ell]}\ \cup\ \big(g_u-g_{u-\delta}\big)_{u\in[\ell]}
\]
are pairwise distinct in $\mathbb Z_m$. The offset $-\delta$ transposes the block and negates $D_\delta$, so it imposes no further condition, and the identity $(H_ZH_Z^T)_{i,i'} = \sum_u\big(G^T_uG_{u+\delta} + F^T_uF_{u+\delta}\big)$ shows the $Z$ sector is governed by the same set of differences.

Choose the $L$ bottom exponents one at a time, in the order $f_0,\dots,f_{\ell-1},g_0,\dots,g_{\ell-1}$. Each element of each $D_\delta$ is a difference of two exponents and so becomes determined exactly when the later of the two is chosen; a step therefore determines at most two elements of $D_\delta$ --- the difference $f_t - f_{t-\delta}$ reaching backwards, and the wrap-around difference $f_{t+\delta-\ell}-f_t$ when $t\geq\ell-\delta$ --- and each newly determined element is an affine function of the current exponent with coefficient $\pm1$. Requiring it to differ from the at most $L-1$ other elements of its own class forbids at most $L-1$ values of that exponent. Summing over the two elements and the $J-1$ offsets, at most $2(J-1)(L-1) < 2(J-1)L$ values are forbidden at any step, so a valid choice exists whenever $m> 2(J-1)L$. Iterating over all $L$ exponents produces the required $\mathcal Q^*$.
\end{proof}


\subsection{Bounds on ZX-Dual Codes}\label{ssec:zx-bounds}
Imposing a ZX-duality can lead to low girth or distance, so we need to be careful to avoid these pitfalls. In particular, the reflection based sector involutions ends up being the most problematic. We show that ZX-dual codes constructed using $r = r_2$ of~\eqref{eq:zx-sector-maps} force girth $4$ whenever $J\geq2$, that $r = r_1r_2$ does the same in the balanced case $J = \tfrac L4$, and that $r = r_1r_2$ at $J=\tfrac L4$ additionally forces $d\leq w$:

\begin{lemma}[Centrality of $t$]\label{lem:t-central}
Let $\mathcal Q\in \mathrm{GALA}_{L,J}(H_k\star C_m)$ have a ZX-duality $\tau =\iota \otimes t$ where $r\in \{r_2, r_1r_2\}$ as defined in~\eqref{eq:zx-sector-maps}.  Then, $t$ commutes with every generator; that is, $t\in Z_{H_k\star C_m}(\mathcal{F}\cup\mathcal{G})$.
\end{lemma}
\begin{proof}
Take first $r = r_2$, for which $\iota = e$ by Proposition~\ref{def:dual-gala}. For $i\in[J]$ and $j\in[\frac L2]$, Definition~\ref{def:kasai} gives
\begin{align*}
(H_X\cdot(e\otimes t))_{i,j} &=F_{j-i}\,t, \\
(H_Z)_{i,j} = G^T_{i-j} = \left(F^T_{j-i}t\right)^T &= t\,F_{j-i}.
\end{align*}
So $H_X\tau = H_Z$ holds only when $F_{j-i}t = tF_{j-i}$ for all $i,j$. It follows that $t$ is central in $\mathcal{F}$, and similarly for the $G$-halves and for $r = r_1r_2$.
\end{proof}

As we will show, the girth, distance, and rate of self-dual codes are all connected to each other via $\Phi_r$:
\begin{equation}\label{eq:zx-gram-block}
    \Phi_r:= \sum_{u\in[\frac{L}{2}]}\left(F_uF^T_{u-r} + F^T_{u-r}F_u\right),
    \qquad r\in[\tfrac{L}{2}].
\end{equation}

\begin{lemma}\label{lem:zx-gram}
Let $\mathcal Q\in \mathrm{GALA}_{L,J}(H_k\star C_m)$ have a ZX-duality $\tau =\iota \otimes t$ with $r: i\mapsto \beta -i$, and let $i,i'\in[J]$ with $\delta = i-i'$. Then, we have
\begin{equation}\label{eq:zx-gram-refl}
    \left(H_XH_X^T\right)_{i,i'} =\Phi_\delta;
\end{equation}
and 
\begin{align}\label{eq:phi-zero}
   \Phi_0 &= L\cdot I,\\
    \Phi_\delta &= \Psi_{\delta + \beta}t \pmod 2.
\end{align}
\end{lemma}


\begin{proof}
The relation $F_iG_{r(i)} = t$ gives $G_i = F^T_{\beta-i}\cdot t$.  We have:
\begin{align}
    \left(H_XH_X^T\right)_{i,i'} 
    &= \sum_{u} F_uF^T_{u-\delta} + \sum_{u} G_uG^T_{u-\delta} \\
    &= \sum_{u} F_uF^T_{u-\delta} +\sum_{u} F^T_{\beta-u}F_{\beta-u+\delta}\\
    &=\sum_{u} F_uF^T_{u-\delta} + \sum_{u} F^T_{u-\delta}F_{u} \label{eq:lemma13}\\
    &=\Phi_\delta,
\end{align}
where we get line~\eqref{eq:lemma13} by re-indexing the second sum with $v = \beta-u+\delta$, so that $\beta - u = v-\delta$. Similarly, we have $(H_XH_Z^T)_{i,i'} = \sum_{j}F_{i-j}G_{j-i'} + \sum_{j}G_{i-j}F_{j-i'}$ gives $\sum_u F_uG_{\delta-u} + \sum_u G_uF_{\delta-u}$. Using $G_{\delta-u} = F^T_{u-(\delta-\beta)}t$ and $G_u = F^T_{\beta-u}t$ together with Lemma~\ref{lem:t-central}, we have:
\begin{equation*}
    \left(H_XH_Z^T\right)_{i,i'} = \sum_u F_uF^T_{u-(\delta-\beta)}\,t + \sum_{u} F^T_{\beta-u}\,t\,F_{\delta-u}
    = \left(\sum_u F_uF^T_{u-(\delta-\beta)} + \sum_{u} F^T_{u-(\delta-\beta)}F_u\right)t = \Phi_{(\delta-\beta)}t,
\end{equation*}
and it follows that $\Psi_\delta =\left(H_XH_Z^T\right)_{i,i'} = \Phi_{\delta - \beta} t\implies \Phi_{\delta} = \Psi_{\delta+\beta}t$, as desired. Finally, for monomial DPG $F_uF^T_u = F^T_uF_u = I$, so
\begin{equation*}
    \Phi_0 = \sum_{u\in[\frac L2]}\left(F_uF^T_u + F^T_uF_u\right) = \sum_{u\in[\frac L2]}2I = L\cdot I.
\end{equation*}
\end{proof}

For reflection-based ZX-dualities, $\Phi_r$ now predicts orthogonality exactly and depends only on $\mathcal F$.

\begin{corollary}[Distance Failures for Reflections]\label{cor:zx-distance}
Let $\mathcal Q\in \mathrm{GALA}_{L,\frac L4}(H_k\star C_m)$ have a ZX-duality $\tau =\iota \otimes t$ where $r= r_1r_2$ as defined in~\eqref{eq:zx-sector-maps}. Then the parent matrices are fully CSS-orthogonal, $\hat H_X\hat H_Z^T = 0$, and, under the proviso of Fact~\ref{fact:augmentation}, $d_\mathcal{Q}\leq w = L$.
\end{corollary}
\begin{proof}
Here $\beta = \frac L4 = J$, so \eqref{eq:zx-gram-refl} at $\delta = 0$ together with \eqref{eq:phi-zero} gives $\Psi_J\, t = \Psi_{0+\beta}\,t = \Phi_0 = wI$, that is $\Psi_J = w\,t \equiv 0 \bmod 2$, since $w = L$ is even. Active orthogonality already supplies $\Psi_\delta = 0$ for every $\delta \in D := \{-(J-1),\dots,J-1\}$, and $D\cup\{J\} = \mathbb Z_{\frac L2}$ because $\frac L2 = 2J$; hence $\Psi_r = 0$ for \emph{every} offset $r$ and $\hat H_X\hat H_Z^T = 0$. So the barrier cannot be broken: every row of $\hat H_X$ is a weight-$L$ operator commuting with all $Z$-checks, so any latent row outside $\rs{H_X}$ is a logical representative of weight $L$ and $d\leq w$ by Fact~\ref{fact:augmentation}. As in Proposition~\ref{prop:distance-bound}, the proviso can fail only if every latent row is already in the active row space, in which case $\mathcal Q$ is its own fully orthogonal parent.
\end{proof}


\begin{proposition}[Girth-4 Dual Codes]\label{prop:zx-girth4}
Let $\mathcal Q\in \mathrm{GALA}_{L,J}(H_k\star C_m)$ with $J\geq 2$ have a ZX-duality $\tau =\iota \otimes t$ where $r$ is a reflection, i.e.\ either $r = r_2$, or $r = r_1r_2$ with $J = \frac L4$, as defined in~\eqref{eq:zx-sector-maps}. Then $\gth{\mathcal{Q}} = 4$, and the number of four-cycles in the Tanner graph of $H_X$ is at least
\begin{equation}\label{eq:zx-four-cycle-count}
    \binom{J}{2}\cdot \frac L2\cdot mk = \frac{J(J-1)L\,mk}{4},
\end{equation}
which for $J = \frac L4$ equals $J^2(J-1)mk$.  The same bound holds for $H_Z$.
\end{proposition}
\begin{proof}
By Lemma~\ref{lemma:girth}, we have $t_4 = \sum_{a<b}\binom{N_{ab}}{2}$ with $N_{ab}$ the integer check-check overlap. For a fixed unordered pair $\{i,i'\}$ of distinct active block rows with $\delta = i - i'$, the pairs with one check in each are indexed bijectively by $(x,y)\in[km]^2$ with overlap $\Phi_\delta[x,y]$. Every row of $\Phi_\delta$ sums to $L$, since by~\eqref{eq:zx-gram-block} it is a sum of $L$ permutation matrices, so $\sum_{x,y}\Phi_\delta[x,y] = L\,mk$. Since $\binom{N}{2}\geq \frac N2$ for every even $N\geq 0$, and $\Phi_\delta$ has even entries because $\Phi_\delta \equiv \Psi_{\delta+\beta}t\equiv 0 \bmod 2$ by~\eqref{eq:phi-zero} and active orthogonality,
\begin{equation*}
    \sum_{x,y}\binom{\Phi_\delta[x,y]}{2}\geq \frac12\sum_{x,y}\Phi_\delta[x,y] =\frac L2mk.
\end{equation*}
Summing over the $\binom J2$ pairs gives \eqref{eq:zx-four-cycle-count}; for $J = \frac L4$ we have $\frac L2 = 2J$ and the bound reads $\frac{J(J-1)}2\cdot 2J\cdot mk = J^2(J-1)mk$. In particular $t_4>0$ as soon as $J\geq2$, so $\gth{\mathcal Q} = 4$. Finally $\tau$ is a permutation with $H_X\tau = H_Z$, so it carries the Tanner graph of $H_X$ isomorphically onto that of $H_Z$ and the same count holds on the $Z$ side.
\end{proof}

The two translation-type ZX dualities escape because they misalign the two index pairings.  Repeating the computation of Lemma~\ref{lem:zx-gram} for $r(i) = i+\beta$, where $G_v = F^T_{v-\beta}\cdot t$, gives instead
\begin{equation}\label{eq:zx-translation}
    \left(H_XH_X^T\right)_{i,i'} = \sum_u \left(F_uF^T_{u-\delta} + F^T_uF_{u-\delta}\right),
    \qquad
    \left(H_XH_Z^T\right)_{i,i'} = \left(\sum_u \left(F_uF^T_{(\delta-\beta)-u} + F^T_uF_{(\delta-\beta)-u}\right)\right)t .
\end{equation}
The two sums are no longer forced to agree term by term. Within the range covered by Proposition~\ref{prop:zx-girth4} --- $J\geq2$, and $J=\tfrac L4$ for $r_1r_2$ --- this leaves the translations $r\in\{e, r_1\}$ as the only ZX-dualities compatible with $\gth{\mathcal Q}\geq 6$. Indeed, certified dual instances of Table~\ref{tab:frontier-selfdual-good} confirms this, with every girth-$6$ entry carrying a translation.
